% M1_05: Section 6.1, the shareholder paragraph.
% INSERT AFTER: main.tex line 956 ("... The comparison holds the sale rule and the incumbent distribution fixed.")
% Placement: new paragraph between the Table 2 proceeds sentence and the dividend diagnostic that starts at line 958.
% NUM placeholders: \NUM{Pi_acc} (informed investor expected gross profit with price access at \(r_1\), at one trading unit per share of float),
%   \NUM{breakeven_scale} (break-even trade scale in trading units per share of float, for the comparison with an economy without the informed investor).
% OPTIONAL CLAUSE: the sentence with the two \NUM values stands alone. Delete that one sentence if the numerics are withdrawn; the rest of the paragraph still reads correctly.
% The literal 0.257809 already appears in the paper (Section 6.1 dividend diagnostic) and needs no placeholder.
% New \bibitem requested: none.
% New symbols: none.

Target shareholders gain from the mechanism through entry, not
through the price. A shareholder who holds through the sale receives
\(V_T\), not the interim price, and entry raises expected proceeds
with either class because \(t_H>t_0\) and \(t_L>t_0\). At \(r_1\),
price access raises expected proceeds by 0.257809 per share. The
informed investor's expected gross profit at one trading unit per
share of float is \NUM{Pi_acc}, a transfer from liquidity traders,
and initial shareholders gain in total, relative to an economy
without the informed investor, at each trade scale below
\NUM{breakeven_scale} trading units per share of float. Holders who
do not trade keep the full proceeds gain. Holders who must trade
bear the transfer and can lose even when initial shareholders gain
in total. The gain requires the recruiting region of
\hyperref[result-proposition-2]{Proposition 2}, in which high-cost
preparation does not pay at the prior. If high-cost preparation
already pays at the prior, an informative price deters some
preparation and lowers expected proceeds.
